% =====================================================================
%  Section 5.6: Numerical Results  --  REVISED for the M-1 grid
%  n = 60 per alpha, n = 30 per (alpha, topic). 540 runs.
%  Replaces the T=35 / 180-run results entirely.
%  ALL CONTENT IN THIS FILE IS NEW -> rendered in green.
% =====================================================================
\begingroup\color{revgreen}

\subsection{Numerical Results}
\label{sec:results}

We first establish that the simulations converge (Section~\ref{sec:res-conv}), since
every subsequent reading is taken at each run's converged state. We then report the
final-state comparison (Section~\ref{sec:res-table}) and organize the findings around the
research questions: efficiency and PoA (RQ1, Section~\ref{sec:res-rq1}), echo chamber
formation (RQ2, Section~\ref{sec:res-rq2}), and the existence of an optimal mixture
(RQ3, Section~\ref{sec:res-rq3}). Section~\ref{sec:res-topology} reports the manipulation
of initial topology.

\subsubsection{Convergence}
\label{sec:res-conv}

\begin{table}[t]
\centering
\caption{Convergence by mixing ratio, pooled over three topologies and two topics
($n = 60$ per row). $t_{\mathrm{conv}}$ is the detected attractor-entry step;
$T$ is total steps run ($= t_{\mathrm{conv}} + 10$ when converged). No run at any
$\alpha$ exceeded $T_{\max} = 120$ more than once.}
\label{tab:convergence}
\begin{tabular}{lccccc}
\toprule
$\alpha$ & $t_{\mathrm{conv}}$ mean & [range] & $T$ mean & [range] & hit $T_{\max}$ \\
\midrule
0.0   & 52.2 & [16--115] & 62.1 & [26--120] & 0/60 \\
0.125 & 45.5 & [15--112] & 55.5 & [25--120] & 0/60 \\
0.25  & 43.8 & [19--108] & 54.9 & [29--120] & 1/60 \\
0.375 & 37.6 & [9--75]   & 48.8 & [19--120] & 1/60 \\
0.5   & 34.5 & [11--79]  & 44.5 & [21--89]  & 0/60 \\
0.625 & 31.2 & [13--46]  & 41.2 & [23--56]  & 0/60 \\
0.75  & 24.0 & [7--47]   & 34.0 & [17--57]  & 0/60 \\
0.875 & 21.7 & [7--38]   & 31.7 & [17--48]  & 0/60 \\
1.0   & 25.1 & [7--31]   & 35.0 & [17--41]  & 0/60 \\
\midrule
all   & 45.3 & --- & --- & --- & 2/540 \\
\bottomrule
\end{tabular}
\end{table}

Table~\ref{tab:convergence} reports settling behavior across the grid. Two results
follow, and both bear directly on the validity of the measurements.

First, \textbf{538 of 540 runs (99.63\%) converged within $T_{\max} = 120$}, the two
exceptions being classified \texttt{attractor} $=$ \texttt{none} and retained in all
aggregates. Convergence is therefore a property of the model across the whole grid rather
than of a favorable subset.

Second, \textbf{a fixed horizon of $T = 35$ is inadequate for most of the grid.} Mean
settling time at $\alpha = 0$ is 52.2 steps and individual runs require up to 115. Only
the region $\alpha \ge 0.75$ settles within 35 steps on average. Any metric read at
$t = 35$ across all $\alpha$ therefore compares converged high-$\alpha$ states against
low-$\alpha$ transients, and the resulting $\alpha$-response curve confounds the effect
of composition with differences in how far each run has progressed. This is the single
most consequential correction to our earlier configuration, and it motivates the dynamic
stopping rule of Section~\ref{sec:stopping}. Equally, $T_{\max} = 120$ is not excessive:
it is required by the low-$\alpha$ tail and is reached by only 2 runs.

Settling time \textbf{decreases monotonically in $\alpha$ from 52.2 to 21.7 steps ---
and then rises again at $\alpha = 1.0$}, to 25.1. A population that is almost entirely
Type-C converges faster than one that is entirely Type-C. Section~\ref{sec:res-rq3}
shows this is the same phenomenon as the collapse in fixed-point attainment at
$\alpha = 1$.

Convergence is also slower for gun control than for abortion at every $\alpha$ (at
$\alpha = 0$: 58.7 vs.\ 45.7 steps). This is consistent with their initial stance
distributions: gun control is near-balanced (54/46) with many agents close to the
decision boundary, whereas abortion has a clear majority (68/32) and settles sooner.

\subsubsection{Final-state comparison}
\label{sec:res-table}

\begin{table}[t]
\centering
\caption{Metrics at each run's converged state ($t_{\mathrm{conv}} + 10$), by topic and
mixing ratio. $n = 30$ per row (3 topologies $\times$ 10 seeds). Mean $\pm$ 95\% CI.
$\downarrow$ denotes lower-is-better. $Q_{\mathrm{norm}}$ is the degree-preserving
null-corrected modularity (Section~\ref{sec:metrics}) and is the basis for all
substantive claims; bare $Q$ is retained for comparability with prior work.}
\label{tab:main}
\begin{tabular}{llcccc}
\toprule
Topic & $\alpha$ & $P_z$ & $Q$ & $Q_{\mathrm{norm}}$ & PoA $\downarrow$ \\
\midrule
\multirow{9}{*}{Gun Control}
 & 0     & $0.261 \pm 0.013$ & $0.486 \pm 0.009$ & $0.241 \pm 0.009$ & $5.681 \pm 0.340$ \\
 & 0.125 & $0.312 \pm 0.016$ & $0.497 \pm 0.009$ & $0.250 \pm 0.009$ & $5.480 \pm 0.312$ \\
 & 0.25  & $0.364 \pm 0.013$ & $0.510 \pm 0.010$ & $0.259 \pm 0.009$ & $5.642 \pm 0.396$ \\
 & 0.375 & $0.409 \pm 0.014$ & $0.518 \pm 0.010$ & $0.263 \pm 0.009$ & $5.477 \pm 0.325$ \\
 & 0.5   & $0.451 \pm 0.012$ & $0.540 \pm 0.009$ & $0.277 \pm 0.007$ & $5.374 \pm 0.395$ \\
 & 0.625 & $0.494 \pm 0.011$ & $0.561 \pm 0.011$ & $0.286 \pm 0.010$ & $4.424 \pm 0.344$ \\
 & 0.75  & $0.538 \pm 0.008$ & $0.608 \pm 0.009$ & $0.314 \pm 0.008$ & $3.821 \pm 0.274$ \\
 & 0.875 & $0.582 \pm 0.006$ & $0.664 \pm 0.009$ & $0.350 \pm 0.009$ & $3.220 \pm 0.222$ \\
 & 1     & $0.619 \pm 0.000$ & $0.763 \pm 0.001$ & $0.427 \pm 0.002$ & $\mathbf{1.149 \pm 0.008}$ \\
\midrule
\multirow{9}{*}{Abortion}
 & 0     & $0.347 \pm 0.017$ & $0.423 \pm 0.007$ & $0.194 \pm 0.006$ & $5.436 \pm 0.535$ \\
 & 0.125 & $0.364 \pm 0.011$ & $0.425 \pm 0.007$ & $0.199 \pm 0.007$ & $6.322 \pm 0.389$ \\
 & 0.25  & $0.426 \pm 0.015$ & $0.423 \pm 0.008$ & $0.191 \pm 0.007$ & $6.077 \pm 0.367$ \\
 & 0.375 & $0.463 \pm 0.016$ & $0.495 \pm 0.015$ & $0.241 \pm 0.013$ & $5.815 \pm 0.385$ \\
 & 0.5   & $0.480 \pm 0.015$ & $0.553 \pm 0.012$ & $0.283 \pm 0.010$ & $5.117 \pm 0.298$ \\
 & 0.625 & $0.510 \pm 0.013$ & $0.581 \pm 0.010$ & $0.297 \pm 0.008$ & $4.670 \pm 0.296$ \\
 & 0.75  & $0.534 \pm 0.010$ & $0.613 \pm 0.009$ & $0.317 \pm 0.008$ & $3.853 \pm 0.288$ \\
 & 0.875 & $0.554 \pm 0.008$ & $0.657 \pm 0.008$ & $0.343 \pm 0.006$ & $2.758 \pm 0.306$ \\
 & 1     & $0.575 \pm 0.000$ & $0.764 \pm 0.000$ & $0.419 \pm 0.002$ & $\mathbf{1.129 \pm 0.001}$ \\
\bottomrule
\end{tabular}
\end{table}

Table~\ref{tab:main} is the primary result. Before interpreting it we note a validation
that constrains how the differences from our earlier configuration may be read.

\paragraph{The numerical baseline reproduces exactly.}
At $\alpha = 1$ the population contains no LLM agents and the dynamics reduce to the FJ
closed form. The values there are unchanged from our earlier configuration:
$P_z = 0.619$ / $0.575$ and $Q = 0.763$ / $0.764$ for gun control and abortion, against
$0.619$ / $0.575$ and $0.764$ / $0.764$ previously. \textbf{The numerical path is
therefore not affected by any of the changes described in Section~\ref{sec:config},
which means every difference at $\alpha < 1$ is attributable to the Type-L agents and
the continuous scale, not to a harness error.}

\subsubsection{RQ1: efficiency and the Price of Anarchy}
\label{sec:res-rq1}

\paragraph{Type-C populations land on the game-theoretic prediction.}
At $\alpha = 1$, PoA is $1.129 \pm 0.001$ (abortion) and $1.149 \pm 0.008$ (gun control),
against the $9/8 = 1.125$ upper bound of \citet{bhawalkar2013coevolutionary}. This is a
stronger statement than reproducing our own earlier numbers. The planner's solution
$z^* = (L_{\mathrm{sym}} + W)^{-1} W s$ is precisely the fixed point of the FJ update
rule, so Type-C agents converge by construction to a neighborhood of the social optimum,
and the residual gap is exactly the Nash--optimum discrepancy of the opinion formation
game: individual and social cost count each edge a different number of times. The
measured value falls essentially on the theoretical one. (Our model uses directed K-NN
edges counted once each, so $9/8$ is not strictly applicable; the agreement is
nonetheless close enough that we report it as corroboration rather than coincidence.)

\paragraph{Text-based reasoning is five times less efficient.}
Pooled across topologies and topics, PoA rises from $1.139 \pm 0.005$ at $\alpha = 1$ to
$5.558 \pm 0.309$ at $\alpha = 0$ --- a factor of 4.9. The full curve and the
adjacent-$\alpha$ tests are:

\begin{center}
\begin{tabular}{lccccccccc}
\toprule
$\alpha$ & 0.0 & 0.125 & 0.25 & 0.375 & 0.5 & 0.625 & 0.75 & 0.875 & 1.0 \\
\midrule
PoA      & 5.558 & 5.901 & 5.860 & 5.646 & 5.246 & 4.547 & 3.837 & 2.989 & 1.139 \\
$\pm$ CI & 0.309 & 0.265 & 0.268 & 0.249 & 0.242 & 0.222 & 0.193 & 0.193 & 0.005 \\
\bottomrule
\end{tabular}
\end{center}

\begin{center}
\begin{tabular}{lcc}
\toprule
Comparison & change & $q$ (Welch, BH-FDR) \\
\midrule
$0 \to 0.125$      & $5.558 \to 5.901$ & 0.127 \\
$0.125 \to 0.25$   & $5.901 \to 5.860$ & 0.826 \\
$0.25 \to 0.375$   & $5.860 \to 5.646$ & 0.280 \\
$0.375 \to 0.5$    & $5.646 \to 5.246$ & $\mathbf{0.036}$ \\
$0.5 \to 0.625$    & $5.246 \to 4.547$ & $\mathbf{8.5\times10^{-5}}$ \\
$0.625 \to 0.75$   & $4.547 \to 3.837$ & $\mathbf{1.1\times10^{-5}}$ \\
$0.75 \to 0.875$   & $3.837 \to 2.989$ & $\mathbf{3.0\times10^{-8}}$ \\
$0.875 \to 1.0$    & $2.989 \to 1.139$ & $\mathbf{3.9\times10^{-26}}$ \\
\bottomrule
\end{tabular}
\end{center}

\textbf{The curve has two regimes.} For $\alpha \le 0.375$ it is a plateau: no
adjacent-$\alpha$ comparison is significant, and PoA stays between 5.6 and 5.9. From
$\alpha = 0.375$ onward it declines monotonically, with every adjacent comparison
significant after FDR correction. Replacing LLM agents with cost-minimizing agents buys
nothing until they exceed roughly a third of the population.

\paragraph{Why this is roughly twice our earlier estimate.}
Our earlier configuration reported PoA $= 2.807$ at $\alpha = 0$. That figure was
produced with Type-L beliefs quantized to five levels
$\{-1, -0.5, 0, 0.5, 1\}$. \emph{Quantization truncates the distance an agent can
travel from the continuous optimum}, so the measured efficiency loss was mechanically
compressed. The present value is not an inflation introduced by new machinery; it is
what the earlier scale was unable to express.

\paragraph{Decomposition: inefficiency is mostly conformity, not disagreement.}
Table~\ref{tab:decomp} separates the social cost into its disagreement and conformity
terms (Section~\ref{sec:metrics}).

\begin{table}[t]
\centering
\caption{Decomposition of social cost at the converged state, pooled ($n = 60$ per row).
``Disagreement'' is $\sum_i\sum_{j \in \mathcal{N}(i)}(z_i - z_j)^2$; ``conformity'' is
$\sum_i \rho_i K (z_i - s_i)^2$, the cost of having moved away from one's own intrinsic
position.}
\label{tab:decomp}
\begin{tabular}{lcccc}
\toprule
$\alpha$ & PoA & Disagreement & Conformity & Conformity share \\
\midrule
0.0   & 5.558 & 3.136 & 2.422 & \textbf{43.6\%} \\
0.125 & 5.901 & 3.690 & 2.211 & 37.5\% \\
0.25  & 5.860 & 3.882 & 1.978 & 33.7\% \\
0.375 & 5.646 & 3.822 & 1.824 & 32.3\% \\
0.5   & 5.246 & 3.610 & 1.636 & 31.2\% \\
0.625 & 4.547 & 3.127 & 1.420 & 31.2\% \\
0.75  & 3.837 & 2.678 & 1.159 & 30.2\% \\
0.875 & 2.989 & 2.175 & 0.814 & 27.2\% \\
1.0   & 1.139 & 0.956 & 0.183 & \textbf{16.1\%} \\
\bottomrule
\end{tabular}
\end{table}

Moving from $\alpha = 1$ to $\alpha = 0$, the disagreement term grows by a factor of
$3.3$ ($0.956 \to 3.136$) while \textbf{the conformity term grows by a factor of $13.2$}
($0.183 \to 2.422$). The conformity share rises monotonically from 16.1\% to 43.6\% with
no exception across the nine levels.

This is the substantive mechanism behind the aggregate. \textbf{LLM-driven agents are not
merely more argumentative; they are markedly more persuadable, and abandoning their own
prior stance is the dominant source of their inefficiency.} A Type-C agent is
mathematically prevented from full conformity by the $\rho_i$ term in its objective. A
Type-L agent receives its intrinsic stance as an anchor in the prompt at every step but
can be argued away from it, and in the homogeneous neighborhoods that K-NN rewiring
creates, it repeatedly is.

\subsubsection{RQ2: echo chamber formation}
\label{sec:res-rq2}

\begin{center}
\begin{tabular}{lccccccccc}
\toprule
$\alpha$ & 0.0 & 0.125 & 0.25 & 0.375 & 0.5 & 0.625 & 0.75 & 0.875 & 1.0 \\
\midrule
$P_z$              & 0.304 & 0.338 & 0.395 & 0.436 & 0.466 & 0.502 & 0.536 & 0.568 & 0.597 \\
$Q_{\mathrm{norm}}$& 0.218 & 0.225 & 0.225 & 0.252 & 0.280 & 0.292 & 0.316 & 0.347 & 0.423 \\
\bottomrule
\end{tabular}
\end{center}

Both quantities rise \textbf{monotonically} in $\alpha$, with narrow confidence intervals
(Table~\ref{tab:main}). Polarization roughly doubles from $0.304$ to $0.597$;
null-corrected modularity rises from $0.218$ to $0.423$.

\textbf{This reverses the conclusion of our earlier configuration}, which reported $P_z$
confined to $0.575$--$0.66$ with no consistent trend and concluded that the mixing ratio
has little influence on opinion spread. That flatness was an artifact of the five-point
belief scale: with only five admissible positions, the variance of the population is
bounded below by the spacing of the grid, and the measured $P_z$ was dominated by
quantization rather than by the dynamics. On a continuous scale, $\alpha$ turns out to be
the strongest single determinant of both polarization and community structure that we
measure.

The direction is consistent with the decomposition in Table~\ref{tab:decomp}:
\textbf{LLM populations are less polarized and less modular precisely because they are
more persuadable.} A population that can be argued out of its positions does not hold the
camps that polarization and modularity measure. Type-C agents, held to their intrinsic
opinions by a formal penalty, sustain the separation.

We report $Q_{\mathrm{norm}}$ rather than bare $Q$ throughout. At $\alpha = 1$, bare
$Q = 0.763$ but $\bar{Q}_{\mathrm{rand}} = 0.335$: \textbf{roughly 44\% of the raw
modularity is reproduced by a random graph with the same degree sequence.} The structure
remains strongly non-random ($z_Q \approx 24$--$26$), but the description ``$Q \approx
0.76$ indicates strong echo chambers'' overstates the effect by more than a factor of
two and should be restated as $Q_{\mathrm{norm}} \approx 0.42$--$0.43$.

$Q_{\mathrm{norm}}$ is nearly flat over $\alpha \in [0, 0.25]$ ($0.218$ / $0.225$ /
$0.225$) before rising, mirroring the PoA plateau over the same range.

\subsubsection{RQ3: is there an optimal mixture?}
\label{sec:res-rq3}

For efficiency alone the answer is $\alpha = 1$, which attains PoA $\approx 1.14$; no
interior value approaches it. But the grid reveals a structural property of hybrid
populations that neither pure population exhibits.

\begin{table}[t]
\centering
\caption{Attractor composition by mixing ratio, pooled ($n = 60$ per row).
\texttt{fixed\_point} means the edge set becomes exactly stationary.}
\label{tab:attractor}
\begin{tabular}{lccccc}
\toprule
$\alpha$ & \texttt{fixed\_point} & \texttt{limit\_cycle} & \texttt{plateau} & \texttt{none} & fixed-point share \\
\midrule
0.0   & 11 & 1 & 48 & 0 & \textbf{18\%} \\
0.125 & 21 & 2 & 37 & 0 & 35\% \\
0.25  & 26 & 1 & 32 & 1 & 43\% \\
0.375 & 24 & 1 & 34 & 1 & 40\% \\
0.5   & 32 & 2 & 26 & 0 & 53\% \\
0.625 & 30 & 2 & 28 & 0 & 50\% \\
0.75  & 36 & 5 & 19 & 0 & 60\% \\
0.875 & 37 & 5 & 18 & 0 & \textbf{62\%} \\
1.0   & 7  & 1 & 52 & 0 & \textbf{12\%} \\
\bottomrule
\end{tabular}
\end{table}

\textbf{Mixed populations freeze; pure populations do not.} Both homogeneous
configurations reach an exactly stationary edge set in a minority of runs (18\% at
$\alpha = 0$, 12\% at $\alpha = 1$), whereas the fixed-point share rises to 62\% at
$\alpha = 0.875$ and then collapses. A population containing just six LLM agents out of
fifty is five times more likely to freeze completely than one containing none.

The plateau height of $C^{\mathrm{out}}$ shows the complementary pattern:

\begin{center}
\begin{tabular}{lcc}
\toprule
$\alpha$ & $C_{\mathrm{plateau}}$ & sd \\
\midrule
0.0          & \textbf{0.9129} & 0.0687 \\
0.125--0.5   & $\approx 0.971$ & 0.019--0.027 \\
0.625--0.875 & 0.977--0.983    & 0.013--0.016 \\
1.0          & \textbf{0.9901} & 0.0043 \\
\bottomrule
\end{tabular}
\end{center}

Only $\alpha = 0$ is clearly distinct --- roughly 8.7\% of neighbors replaced per step,
against about 1\% for Type-C --- with the interior essentially flat. So ``Type-L agents
churn their neighborhoods more'' holds, but ``churn decreases monotonically with
$\alpha$'' does not. This corrects the tentative reading from the $n = 3$ pilot
(Section~\ref{sec:stopping}).

We do not claim a verified mechanism for the freezing. A testable conjecture is that a
small number of Type-L agents supplies fixed anchors that interrupt the continual
micro-adjustment of the Type-C population, letting the system land on an exact fixed
point rather than approaching one asymptotically. The rise in $t_{\mathrm{conv}}$ at
$\alpha = 1$ (Table~\ref{tab:convergence}) is the same phenomenon seen from the time axis.

\subsubsection{Initial topology does not determine the outcome}
\label{sec:res-topology}

\begin{center}
\begin{tabular}{lccc}
\toprule
$\alpha$ & scale-free (BA) & random (ER) & small-world (WS) \\
\midrule
0.0  & $5.348 \pm 0.503$ & $5.766 \pm 0.737$ & $5.562 \pm 0.395$ \\
0.25 & $5.944 \pm 0.356$ & $5.961 \pm 0.531$ & $5.673 \pm 0.557$ \\
0.5  & $5.078 \pm 0.373$ & $5.422 \pm 0.417$ & $5.237 \pm 0.521$ \\
0.75 & $3.912 \pm 0.420$ & $4.051 \pm 0.332$ & $3.549 \pm 0.247$ \\
1.0  & $1.138 \pm 0.008$ & $1.134 \pm 0.007$ & $1.144 \pm 0.010$ \\
\bottomrule
\end{tabular}
\end{center}

Across all nine $\alpha$ levels and both PoA and $P_z$, the confidence intervals of the
three topology families overlap and \textbf{no comparison reaches significance}
($n = 20$ per cell).

\paragraph{Absence of a detected difference is not equivalence.}
We tested equivalence directly with TOST at the pre-registered margin $\delta = 0.2$:

\begin{center}
\begin{tabular}{lccccc}
\toprule
$\alpha$ & most separated pair & difference & TOST $p$ & verdict \\
\midrule
0.0  & BA / ER & $-0.418$ & 0.694 & not established \\
0.25 & ER / WS & $+0.288$ & 0.594 & not established \\
0.5  & BA / ER & $-0.344$ & 0.704 & not established \\
0.75 & ER / WS & $+0.501$ & 0.932 & not established \\
1.0  & BA / WS & $\mathbf{-0.006}$ & $\mathbf{<0.0001}$ & \textbf{equivalent} \\
\bottomrule
\end{tabular}
\end{center}

\textbf{Only $\alpha = 1$ passes.} At low $\alpha$ the run-to-run variance of PoA
(sd $\approx 1.0$--$1.5$) is large enough that the data cannot exclude a true difference
of up to roughly $0.5$ between topology families. We state the conservative and the
mechanistic readings separately.

\emph{Conservative claim, fully supported by the data.} At $\alpha = 1$ the three initial
topologies are statistically equivalent (TOST, $\delta = 0.2$, $p < 0.0001$). At every
other $\alpha$, no significant difference is detected, but the sample size does not
support a formal equivalence claim.

\emph{Mechanistic reading.} The rewiring rule deletes all edges at every step and rebuilds
them from opinion distances, so the initial topology is overwritten after the first
rewiring. Mean degrees were matched to within 4.17\% while clustering coefficients
remained clearly distinct (0.179 / 0.079 / 0.390), and no difference was detected under
that design --- \textbf{echo chamber formation here does not require a pre-existing hub
structure}. That this also holds at $\alpha = 0$, where no agent follows the FJ rule,
shows it is a property of the rewiring mechanism rather than of the FJ closed form.

We flag a scale issue with the margin. $\delta = 0.2$ is a 17.5\% relative bound at
$\alpha = 1$ but only 3.6\% at $\alpha = 0$. Asserting equivalence at low $\alpha$ would
require either a relative margin (e.g.\ 5\%, i.e.\ $\delta \approx 0.28$) or a larger
$n$; both are decisions that must be fixed in advance rather than chosen after
inspecting the data, and we therefore leave the low-$\alpha$ claim open.

\subsubsection{A topic-specific PoA peak}
\label{sec:res-peak}

Pooling topics, the apparent rise from $\alpha = 0$ to $\alpha = 0.125$ is not
significant ($p = 0.095$). Separating them shows why:

\begin{center}
\begin{tabular}{lccccc}
\toprule
Topic & $\alpha = 0$ & $\alpha = 0.125$ & Welch $t$ & $p$ & Cohen's $d$ \\
\midrule
\textbf{Abortion} & 5.436 & $\mathbf{6.322}$ & $+2.74$ & $\mathbf{0.0084}$ & $\mathbf{0.71}$ \\
Gun control       & 5.681 & 5.480 & $-0.89$ & 0.379 & $-0.23$ \\
Pooled            & 5.558 & 5.901 & $+1.68$ & 0.095 & 0.31 \\
\bottomrule
\end{tabular}
\end{center}

\textbf{Abortion exhibits a significant peak at $\alpha = 0.125$ with a medium-to-large
effect ($d = 0.71$); gun control has no peak at all}, its maximum being at $\alpha = 0$.

This both confirms and relocates the non-monotonicity our earlier configuration reported.
A peak exists, and it is specific to abortion --- but it sits at $\alpha = 0.125$, not
$\alpha = 0.75$, and it is about twice as large.

A testable mechanism: abortion has the more asymmetric stance distribution (68/32 against
gun control's 54/46). On an issue with a clear majority, introducing a small number of
Type-C agents disrupts the otherwise consistent coordination among Type-L agents and
drives social cost up; on a near-balanced issue the effect is absent. Because both topics
use identical personas and stubbornness values (Section~\ref{sec:dataset}), this
difference is attributable to the stance distribution and not to population composition.

\subsubsection{Key findings}
\label{sec:res-findings}

\paragraph{F1 --- Echo chamber structure is produced by the rewiring rule, but its
magnitude is set by agent type.}
No initial topology among BA, ER and WS produces a detectably different outcome at any
$\alpha$, and at $\alpha = 1$ the three are formally equivalent: homophilic rewiring
alone suffices, and no pre-existing hub structure is required. \emph{However, the claim
that agent type is irrelevant to structure does not survive continuous measurement.}
Polarization roughly doubles across the $\alpha$ range ($0.304 \to 0.597$) and
null-corrected modularity nearly doubles ($0.218 \to 0.423$), both monotonically. The
connection rule determines \emph{that} echo chambers form; the agent population
determines \emph{how strong} they become.

\paragraph{F2 --- Language-based reasoning operates as a conformity mechanism, and this
is now measured rather than inferred.}
LLM populations incur PoA $= 5.56$ against $1.14$ for numerical populations. The
decomposition locates the loss: 43.6\% of the social cost at $\alpha = 0$ comes from
agents having moved away from their own intrinsic position, against 16.1\% at
$\alpha = 1$ --- a 13.2-fold difference in the conformity term. The capacity for
natural-language reasoning that makes Type-L agents expressive is the same capacity that
makes them yield under sustained directional pressure.

\paragraph{F3 --- The efficiency response to $\alpha$ has a plateau and a decline, and
the peak is topic-specific.}
PoA is statistically flat for $\alpha \le 0.375$ and declines significantly at every step
thereafter. Within the plateau, abortion shows a significant local peak at
$\alpha = 0.125$ ($d = 0.71$) that gun control does not. \emph{This supersedes our earlier
characterization of a norm contention cost peaking in the mid-to-high range
($\alpha \approx 0.5$--$0.75$) consistently across topics.} That reading was obtained at a
fixed $T = 35$, which for $\alpha \le 0.5$ measures unconverged transients
(Table~\ref{tab:convergence}); the effect does not survive measurement at convergence.

\paragraph{F4 --- Hybrid populations solidify more readily than either pure population.}
Exact fixed points occur in 18\% of runs at $\alpha = 0$ and 12\% at $\alpha = 1$, but in
62\% at $\alpha = 0.875$. The mechanism is not established; we report it as a robust
regularity over 540 runs and offer the anchoring conjecture of
Section~\ref{sec:res-rq3} as a hypothesis.

\paragraph{F5 --- Numerical agents land on the game-theoretic bound.}
At $\alpha = 1$, PoA $= 1.129$--$1.149$ against the $9/8 = 1.125$ bound of
\citet{bhawalkar2013coevolutionary}. The simulation does not merely reproduce our own
previous numbers; it lands where opinion formation game theory predicts, which
constrains how far the low-$\alpha$ results can be attributed to implementation error.

\endgroup
